\clearpage

\setcounter{section}{7}

\subsubsection{Beitrag der Elektronen zur Wärmekapazität des Festkörpers}\,\\ 
Für die spezifische Wärme gilt
$$
C_{V} = \left( \frac{\mathrm{d}Q}{\mathrm{d}T} \right) V = \frac{\mathrm{d}U}{\mathrm{T}}
$$ 
mit der inneren Energie pro Volumen
$$
U_0 = N \overline{E}.
$$ 
Es folgt
\begin{equation}
	\label{eq:II.7.17}
	C_{V} = \frac{\mathrm{d}U_0}{\mathrm{d}T} = N \frac{\mathrm{d}}{\mathrm{d}T} \overline{E} \overset{\ref{eq:II.7.16}}{=} N k_{\mathrm{B}} \frac{\pi^2}{2} \frac{T}{T_{\mathrm{F}}}
\end{equation}
was dem quantenmechnischem Ergebnis entspricht. Für das klassische Ergebnis gilt hingegen
$$
U_0 = N \frac{3}{2} k_{\mathrm{B}}T \quad \to \quad C_{V}^{\text{Kl.}} = \frac{3}{2} N k_{\mathrm{B}}
$$ 
wobei der direkte Vergleich
$$
\frac{C_{V}^{\text{QM}}}{C_{V}^{\text{Kl.}}} \approx \frac{T}{T_{\mathrm{F}}} = \frac{\num{300} \, \text{(RT)}}{\num{30000}} = \frac{1}{100} 
$$ 
liefert. \\
\textcolor{red}{ref zu abbildung aus vorheriger vl}\\
Nur ein 'paar' Elektronen gewinnen die Energie $k_{\mathrm{B}}T$ wegen dem Pauli-Prinzip. Es müssen freie Zustände verfügbar sein, dies ist nur für Elektronen an der Oberfläche der Fermi-Kugel (Fermi-Fläche) der Fall. \\
Metalle bei niedrigen Temperaturen
\begin{equation}
	\label{eq:II.7.20}
	C_{V} = \underbrace{\gamma T}_{\mathrm{e}^{-}} + \underbrace{\alpha T^3}_{\text{Phononen}}
\end{equation}


\paragraph{Beitrag der Elektronen zur Wärmeleitfähigkeit}\,\\ 
\begin{equation}
	\label{eq:II.7.21}
	K = \frac{1}{3} \overline{v} \lambda C_{V}
\end{equation}
mit der mittleren Geschwindigkeit $\overline{v}$,  der freien Weglänge zwischen den Stößen $\lambda$ und der spezifischen Wärme pro Volumen $C_{V}$. Es gilt $\overline{v} = v_{\mathrm{F}}$, da nur Elektronen an der Fermigrenze zu $C_{V}$ beitragen $\text{const.}(T) = \SI{e 8}{ \centi \meter \per \second}$.\textcolor{red}{ref zu vorheriger formel einbinden}\\
Unter Berücksichtigung der Definition der Fermi-Energie
$$
E_{\mathrm{F}} = k_{\mathrm{B}}T_{\mathrm{F}} = \frac{1}{2} m v_{\mathrm{F}}^2
$$ 
folgt für die Fermi-Temperatur
\begin{equation}
	\label{eq:II.7.22}
	T_{\mathrm{F}} = \frac{m v_{\mathrm{F}}^2}{2k_{\mathrm{B}}}
\end{equation}
und es ergibt sich 
\begin{equation}
	\label{eq:II.7.23}
	K = \frac{\pi^2 N k_{\mathrm{B}}^2}{3mv_{\mathrm{F}}}\lambda T = \frac{\pi^2 N k_{\mathrm{B}}^2}{3m}T \tau 
\end{equation}
mit $\lambda = v_{\mathrm{F}}\tau $ mit der mittleren Stoßzeit $\tau $. Es gilt bspw. für die mittlere freie Weglänge bei zwischen Stößen bei RT $\lambda(\SI{300}{\kelvin}) = \SI{e -5}{ \centi \meter}$.

\paragraph{Vergleiche ($T\ge \Theta$)}\,\\ 
$$
\frac{K_{\text{e}}}{K_{\text{G}}} = \frac{V^{\mathrm{e}} C_{V}^{\text{e}} \lambda^{\text{e}}}{V^{\text{G}} C_{V}^{\text{G}} \lambda^{\text{G}}} \approx \frac{\SI{e 8}{\centi\meter \second ^{-1}} \cdot 1 \cdot \SI{e-5}{\centi\meter}}{\SI{5e 5}{\centi \meter \second ^{-1}} \cdot 100  \cdot  \SI{3 e-7}{\centi \meter}} \approx 67
$$ 
Die thermische Leitfähigkeit von Metallen ist um den Faktor 10 bis 100 größer als die der Isolatoren
\begin{table}[htpb]
	\centering
	\caption{caption}
	\label{tab:label}
	\begin{tabular}{|c|c|}
		\hline
	Material & $K$ in $\unit{\watt \centi\meter ^{-1} \kelvin ^{-1}}$ \\
	\hline
	Metall & 1 bis 10 \\
	\hline
	Quarz & \num{0,13}\\
	\hline
	\ce{NaCl} & \num{0,07}\\
	\hline
	\makecell{2-atomiges \\ Gas (\ce{N2}, \ce{O2}) & \num{1,7 e-4}\\
	\hline
	\end{tabular}
\end{table}

\subsubsection{Elektrische Leitfähigkeit (Ohm'sches Gesetz)}
\paragraph{a) Druck Modell (kinetische Gastheorie)}\,\\ 
Hierbei handelt es sich um das klassische Modell
\begin{enumerate}
	\item für $| \Vec{E} | = 0$ ergibt sicht für die Driftgeschwindigkeit $v_{\text{D}} =0 $
	\item für $| \Vec{E} | \neq 0$  werden die Elektronen im elektrischen Feld $\Vec{E}$ beschleunigt mit $|\Vec{F}| = -\mathrm{e} | \Vec{E} | $ 
\end{enumerate}

Wir betrachten im Folgenden den eindimensionalen Fall. Es herrscht eine Bewegungshemmende Reibungskraft mit
$$
F_{\text{R}} = m \frac{v_{\mathrm{D}}}{\tau }
$$ 
mit der mittleren Stoßzeit oder Relaxionszeit $\tau $ und der mittleren vom elektrischen Feld zusätzlich bewirkten Geschwindigkeitskomponente $v_{\mathrm{D}}= v- v_{\text{th}}$, wobei $v_{\text{th}}$ der mittleren thermischen Geschwindigkeit entspricht. Die Elektronenbewegung wird beschrieben von der DGL
$$
m\left( \frac{\mathrm{d}v}{\mathrm{d}t} + \frac{1}{\tau } v_{\mathrm{D}} \right) = - \mathrm{e} E
$$ 
wobei im stationären Fall
$$
\dot{v} = \frac{\mathrm{d}v}{\mathrm{d}t}= 0
$$ 
gilt. Aus 
$$
v_{\mathrm{D}} = - \frac{\mathrm{e}E}{m} \tau 
$$ 
folgt der Strom
 $$
j = - N \mathrm{e} v_{\mathrm{D}} = N \frac{\mathrm{e}^2 E}{m} \tau .
$$ 
Wir führen die Leitfähigkeit $\sigma$ mit $j = \sigma E$ und
\begin{equation}
	\label{eq:II.7.25}
	\sigma = \frac{N \mathrm{e}^2 \tau }{m} \quad \mu = \frac{\mathrm{e}}{m} \tau 
\end{equation}
ein. 

\paragraph{b) Modell mit Fermi-Kugel (Sommerfeld)}\,\\ 
Die Fermi-Kugel schließt im $\Vec{k}$-Raum die besetzten Zustände des Elektronengases ein. Der Gesamtimpuls ist Null, oder anders: zu jedem $\Vec{k}$ gibt es ein $-\Vec{k}$

\begin{figure}[H]
    \centering
    \def\svgwidth{.377\columnwidth}
    \import{figures/}{vl26_abbildung_1.pdf_tex}
    \label{fig:vl26_abbildung_1}
\end{figure}

\begin{figure}[H]
    \centering
    \def\svgwidth{.656\columnwidth}
    \import{figures/}{vl26_abbildung_2.pdf_tex}
    \label{fig:vl26_abbildung_2}
\end{figure}

Der Impuls des Elektrons entspricht
$$
m\Vec{v} = \hbar  \Vec{k}
$$ 
mit der Kraft
$$
\begin{aligned}
	\Vec{F} = m \frac{\mathrm{d}\Vec{v}}{\mathrm{d}t} &= \hbar  \frac{\mathrm{d}\Vec{k}}{\mathrm{d}t} = - \mathrm{e} \Vec{E} \\
	\Vec{k}(t) - \Vec{k}(0) &= - \frac{\mathrm{e} \Vec{E} t}{\hbar } \\
\end{aligned}
$$
Die Fermi-Kugel wird als Ganzes verschoben, da jedes Elektron die gleiche Änderung $\delta k$ erfährt. Infolge der Stöße zwischen Elektronen und Verunreinigungen, Gitterfehlern der Phononen stellt sich ein stationärer Zustand ein 
$$
\begin{aligned}
	\delta \Vec{k} &= - \frac{\mathrm{e}\Vec{E} \tau }{\hbar } \\
	v_{\mathrm{D}} &= - \frac{\mathrm{e}\Vec{E} \tau }{m} \\
	\sigma &= \frac{N \mathrm{e}^2 \tau }{m} 
\end{aligned}
$$ 
Es gilt bspw. für reines Kupfer bei \SI{300}{\kelvin}:
$$
\sigma = \SI{6 e 5}{\frac{1}{\ohm \centi \meter}}, \quad \tau = \SI{2 e-14}{\second}, \quad v_{\mathrm{F}} = \SI{1,6 e8}{ \frac{\centi \meter}{\second}}, \quad \lambda = \tau v_{\mathrm{F}} = \SI{3 e-6}{\centi \meter}
$$ 
Es geschehen keine Stöße mit den Gitterbausteinen, da $a = \SI{e -8}{\centi \meter}$ \\
$\SI{4}{\kelvin}$ : $\tau = \SI{2 e-9}{\second}$, $\lambda = \SI{0,2}{\centi \meter}$.
Beachte:
$$
\begin{aligned}
	v_{\mathrm{F}} &\approx \SI{e 8}{ \frac{\centi \meter}{\second}}  \\
	v_{\mathrm{D}} &= \frac{\mathrm{e}E}{m} \tau \quad \left(E = \SI{0,1}{ \frac{\volt}{\centi \meter}}\right)\\
	v_{D} &\simeq \SI{3}{ \frac{\centi \meter}{\second}}
\end{aligned}
$$ 

\subsubsection{Wiedemann-Franz'sches Gesetz}
Das Wiedemann-Franz'sche Gesetz sagt aus, dass für metalle bei nicht zu tiefen Temperaturen ($T>\Theta$) das Verhältnis von thermischer zu elektrischer Leitfähigkeit direkt proportional zur Temperatur ist
$$
\begin{aligned}
	k &= \frac{\pi^2 N k_{\mathrm{B}}^2}{3m} T \tau  \\
	\sigma &= \frac{N\mathrm{e}^2}{m} \tau  
\end{aligned}
$$ 

$$
\frac{k}{\sigma} = \frac{\pi^2}{3}\left( \frac{k_{\mathrm{B}}}{\mathrm{e}} \right) ^2 T = L T
$$ 
wobei
$$
L = \SI{2,45 e-8}{ \frac{\watt \ohm}{\kelvin^2}}
$$ 
der materialunabhängigen Lorentz-Zahl entspricht. Der Zusammenhang gilt nur für $T> \Theta$, wobei für $T\ll\Theta$ thermische und elektrische Relaxationsprozesse nicht identisch sind da unterschiedliche Streuprozesse zum Gleichgewicht führen. $L$ nimmt ab mit kleiner werdender Temperatur.

\paragraph{Temperaturabhängigkeit des elektrischen Widerstands}\,\\ 
Erinnerung: Die mittlere freie Weglänge ist umgekehrt proportional zur mittleren Phononendichte
$$
\lambda \sim \frac{1}{\left<n \right>}:\quad T\ge \Theta \, \to  \, \left<n \right> \sim T
$$ 

\begin{enumerate}
	\item $T\ge \frac{\Theta}{2}$ : $\lambda \sim \frac{1}{T}\, \to  \, \tau \sim \frac{1}{T} \, \to  \, \sigma \sim \frac{1}{T}$
		\begin{equation}
			\label{eq:II.7.28}
			\rho \sim  T
		\end{equation}
	\item $T\le \frac{\Theta}{10}$ 
		$$
		\rho \sim T^{5}
		$$ 
		für reine Substanzen. Kleinwinkelstreuung, simple Überlegung von oben nicht gültig.
\end{enumerate}

Restwiderstand (bie tiefen Temperaturen, unreine Metalle) Streuung auch an Störstellen des Kristalls.
$\to $ Matthiesensche Regel
\begin{equation}
	\label{eq:II.7.29}
	\rho_{\text{Ges}} = \rho_{\text{Stör}} + \rho_{\text{Phonon}}
\end{equation}

\begin{figure}[H]
    \centering
    \def\svgwidth{.328\columnwidth}
    \import{figures/}{vl26_abbildung_3.pdf_tex}
    \label{fig:vl26_abbildung_3}
\end{figure}

$$
\frac{\rho_{\text{Ges}}(\SI{300}{\kelvin})}{\rho_{\text{Ges}}(\SI{4}{\kelvin})} = \frac{\rho_{\text{Ges}}(\SI{300}{\kelvin})}{\rho_{\text{Stör}}}
= \begin{cases}
	\sim 10^{3} - 10^{4}& \text{reine Metall}\\
	\sim 1-2 & \text{Legierungen}
\end{cases}
$$
Metalle mit magnetischen Störstellen (\ce{Fe} in \ce{La})\\
Wechselwirkung $\Vec{\mu}_{\mathrm{e}}$ mit magnetischen Momenten der Störstellen $\Vec{\mu}_{\text{Stör}}$ \\
$\to $ Kondo-Effekt (später)
